\everymath{\displaystyle}

%%%%%%%%%%%%%%%%%%%%%%%%%%%%%%%%%%%%%%%%
%           main body
%%%%%%%%%%%%%%%%%%%%%%%%%%%%%%%%%%%%%%%%

%-----------------------------------
%	TITLE PAGE
%-----------------------------------

\title[第一章\\
集与点集]{\texorpdfstring{\LARGE \bfseries 第一章\quad 集与点集}{第一章 集与点集}} % The short title appears at the bottom of every slide, the full title is only on the title page
\author[\S 1.4\\
开集的构造] % Your institution as it will appear on the bottom of every slide, maybe shorthand to save space
{\Large \bfseries \S 1.4 开集的构造}
% \author{Singular Value Decomposition} % Your name
% \date{\today} % Date can be changed to a custom date
\date{}

%------------------------------------------------

\begin{document}

%------------------------------------------------

\begin{frame}

\titlepage % Print the title page as the first slide

\end{frame}

%------------------------------------------------

% \begin{frame}
% \frametitle{内容提要} % Table of contents slide, comment this block out to remove it
% \tableofcontents % Throughout your presentation, if you choose to use \section{} and \subsection{} commands, these will automatically be printed on this slide as an overview of your presentation
% \end{frame}

%-----------------------------------
%	PRESENTATION SLIDES
%-----------------------------------

%------------------------------------------------

\begin{frame}{绪言}

{
\larger[1]
我们已经学习了开集、闭集的概念 (\textbf{只是描述性的, 非构造性的}), 以及它们的一些性质. 本节我们将给出 $\mathbb{R}$ 中开集的\textbf{构造性}描述, 即任一开集都可以表示为一些开区间的并, 反之任意一些开区间的并也是开集. 这一结果是 $\mathbb{R}$ 的特殊性质 (依赖于 $\mathbb{R}$ 的全序结构); 在高维欧氏空间 $\mathbb{R}^n$ ($n \geqslant 2$) 中, 开集不一定能表示为互不相交开球之并.
}

\end{frame}

%------------------------------------------------

\section[结构定理]{开集的结构定理}

%------------------------------------------------

\begin{frame}{$\mathbb{R}$ 中开集的结构定理}

\begin{thm}[$\mathbb{R}$ 中的开集结构定理]
$\mathbb{R}$ 中任一非空开集 $G$ 均可唯一地表示为至多可列个互不相交的开区间之并:
\[
G = \bigcup_{k} (a_k, b_k), \quad (a_k, b_k) \cap (a_j, b_j) = \emptyset\ (k \neq j),
\]
其中 $a_k, b_k \in \mathbb{R} \cup \{-\infty, +\infty\},$ 称 $(a_k, b_k)$ 为 $G$ 的\alert{构成区间}.\end{thm}

\pause

\startproof[bg=yellow, fg=red]
对每个 $x \in G,$ 令 $I_x = \bigcup\{(a, b) : x \in (a, b) \subset G\}.$ 由 $G$ 开知 $x$ 是 $G$ 的内点, 故 $\{x\}$ 本身可嵌入某开区间 $\subset G,$ 即 $I_x$ 非空. $I_x$ 是包含 $x$ 且含于 $G$ 的最大开区间.

\pause

\textbf{互不相交}: 若 $I_x \cap I_y \neq \emptyset,$ 则 $I_x \cup I_y$ 也是含于 $G$ 的开区间 (两开区间有交则并仍是开区间), 由极大性 $I_x = I_y.$

\pause

\textbf{至多可列}: 各 $I_x$ 互不相交, 每个含至少一个有理数, 不同的含不同的有理数, 故个数至多可列.
\qed
\end{frame}

%------------------------------------------------

\begin{frame}{唯一性与注记}

\startproof[bg=yellow, fg=red](唯一性)
若 $G = \bigcup_k (a_k, b_k)$ 是另一个至多可列互不相交开区间分解, 则对每个 $x \in G,$ $x$ 所在的区间 $(a_k, b_k)$ 是一个包含 $x$ 的开区间且 $\subset G,$ 故 $(a_k, b_k) \subset I_x.$ 反之 $I_x \subset (a_k, b_k)$ 由极大性. 故两种分解的构成区间相同.
\qed
\pause
\vspace{0.5em}

\begin{remark}
此定理是 $\mathbb{R}$ 的特殊性质 (依赖于 $\mathbb{R}$ 的全序结构); 在高维欧氏空间 $\mathbb{R}^n$ ($n \geqslant 2$) 中, 开集不一定能表示为互不相交开球之并.
\end{remark}

\pause

\begin{eg}
\begin{itemize}
\item $G = (0, 1) \cup (2, 5):$ 构成区间为 $(0, 1)$ 和 $(2, 5).$
\item $G = \mathbb{R} \setminus \mathbb{Z} = \bigcup_{n \in \mathbb{Z}} (n, n+1):$ 构成区间为 $(n, n+1)$ ($n \in \mathbb{Z}$), 共可列个.
\item $G = \mathbb{R}:$ 唯一的构成区间是 $(-\infty, +\infty).$
\end{itemize}
\end{eg}

\end{frame}

%------------------------------------------------

\section[Borel 集]{Borel 集类}

%------------------------------------------------

\begin{frame}{$\sigma$-代数}

我们接下来再来看几类重要的集合.

\vspace{0.5em}

之前已经讲过, 开集的任意并、有限交仍是开集; 闭集的任意交、有限并仍是闭集. 一个自然的问题是:

\begin{question}
开集的任意交, 闭集的任意并是什么样的集合? 是否常见? 有无重要性质?
\end{question}

\pause
\vspace{0.5em}

先来引入所谓的 $\sigma$-代数 (sigma-algebra) 的概念, 以便更系统地讨论这些集合.

\begin{Def}[$\sigma$-代数]
设 $X$ 是一个非空集合. $\Gamma$ 是 $X$ 的子集族, 满足:
\begin{itemize}
\item $\emptyset \in \Gamma;$
\item $A \in \Gamma \implies X \setminus A \in \Gamma;$
\item $A_1, A_2, \ldots \in \Gamma \implies \bigcup_{n=1}^\infty A_n \in \Gamma$ (由 De Morgan 律, 也等价于 $\bigcap_{n=1}^\infty A_n \in \Gamma$).
\end{itemize}
则称 $\Gamma$ 是 $X$ 的一个\alert{$\sigma$-代数}, 或称 Borel field.
\end{Def}

\end{frame}

%------------------------------------------------

\begin{frame}{$\sigma$-代数}

\begin{prop}
对非空集合 $X$ 的任意子集族 $\Omega,$ 总存在一个包含它的最小 $\sigma$-代数 $\Gamma(\Omega).$ 也就是说, 对任意 $\sigma$-代数 $\Gamma'$ 满足 $\Omega \subset \Gamma',$ 都有 $\Gamma(\Omega) \subset \Gamma'.$
\end{prop}

\pause

\startproof[bg=yellow, fg=red]
实际上有 $\Gamma(\Omega) = \bigcap\{\Gamma' : \Gamma' \text{ 是 } X \text{ 的 } \sigma\text{-代数},\ \Omega \subset \Gamma'\}.$ 只要验证它是 $\sigma$-代数即可:
\pause
\begin{itemize}
\item $\emptyset \in \Gamma(\Omega)$: 因为 $\emptyset \in \Gamma'$ 对任意 $\Gamma'$ 都成立.
\pause
\item $A \in \Gamma(\Omega) \implies A \in \Gamma'$ 对任意 $\Gamma'$ 都成立, 故 $X \setminus A \in \Gamma'$ 对任意 $\Gamma'$ 都成立, 即 $X \setminus A \in \Gamma(\Omega).$
\pause
\item $A_1, A_2, \ldots \in \Gamma(\Omega) \implies A_n \in \Gamma'$ 对任意 $\Gamma'$ 都成立, 故 $\bigcup_{n=1}^\infty A_n \in \Gamma'$ 对任意 $\Gamma'$ 都成立, 即 $\bigcup_{n=1}^\infty A_n \in \Gamma(\Omega).$
\end{itemize}
\qed

\pause

\begin{Def}
称 $\Gamma(\Omega)$ 是由 $X$ 的子集族 $\Omega$ 生成的 $\sigma$-代数. 也称其为包含 $\Omega$ 的最小 $\sigma$-代数.
\end{Def}

\end{frame}

%------------------------------------------------

\begin{frame}{Borel 集类}

\begin{Def}[Borel 集类]
$\mathbb{R}$ 中由开集生成的 $\sigma$-代数 $\Gamma(\text{开集})$ 称为\alert{Borel 集类}, 其元素称为\alert{Borel 集}.
\end{Def}

\begin{eg}[$G_\delta$ 集与 $F_\sigma$ 集]
\begin{itemize}
\item \alert{$G_\delta$ 集}: 可列个开集之交 ~~ ($G:$ 德文 Gebiet = 开域; $\delta:$ Durchschnitt = 交).
\item \alert{$F_\sigma$ 集}: 可列个闭集之并 ~~ ($F:$ 法文 fermé = 闭; $\sigma:$ Summe = 和).
\end{itemize}
很容易看出来 $G_\delta$ 集与 $F_\sigma$ 集都是 Borel 集, 但它们不一定是开集或闭集.
\end{eg}

\pause

\begin{prop}
每个闭集是 $G_\delta$ 集; 每个开集是 $F_\sigma$ 集.
\end{prop}

\startproof[bg=yellow, fg=red](证明)
设 $F$ 闭. 令 $G_n = \{x : d(x, F) < 1/n\},$ 则各 $G_n$ 是开集, 且
\[
\bigcap_{n=1}^\infty G_n = \{x : d(x,F) = 0\} = \overline{F} = F.
\]
故 $F$ 是 $G_\delta$ 集. 开集情形利用 De Morgan 律即可得.
\qed

\end{frame}

%------------------------------------------------

\begin{frame}{Borel 集类}

进一步迭代可得 $G_{\delta\sigma}$-集, $G_{\delta\sigma\delta}$-集, $\ldots$ 以及 $F_{\sigma\delta}$-集, $F_{\sigma\delta\sigma}$-集, $\ldots$ 这些集合都属于 Borel 集类.
\pause
但是, \alert{并不是所有的 Borel 集都可以通过有限次迭代 $G_\delta$ 和 $F_\sigma$ 的构造得到.}

\pause

\begin{eg}
\begin{itemize}
\item $\mathbb{Q}$ 是 $F_\sigma$ 集 (可列个单点集之并), 但不是 $G_\delta$ 集.
\item $\mathbb{R} \setminus \mathbb{Q}$ (无理数集) 是 $G_\delta$ 集, 但不是 $F_\sigma$ 集.
\item 设 $f(x)$ 是定义在开集上的实值函数, 则 $f(x)$ 的连续点集是 $G_\delta$ 集.
\item 设 $f(x) \in C(\mathbb{R}),$ 则 $f(x)$ 的可微点集是 $F_{\sigma\delta}$ 集.
\end{itemize}
\end{eg}

\pause

\startproof[bg=yellow, fg=red]
简要证明一下最后一个结论. 记 $g(x, h) = \frac{f(x+h) - f(x)}{h},$ 并令
\[
E_{m,n} = \bigcap_{\substack{h, k \in \mathbb{Q} \\ 0 < |h|, |k| < 1/n}} \left\{x : |g(x,h) - g(x,k)| \leqslant \frac{1}{m}\right\}.
\]
则 $E_{m,n}$ 是闭集, 且 $f$ 的可微点集 $E = \bigcap_{m=1}^\infty \bigcup_{n=1}^\infty E_{m,n}$ 是 $F_{\sigma\delta}$ 集.
\qed

\end{frame}

%------------------------------------------------

\section[n 进制]{实数的 $n$ 进制表示}

%------------------------------------------------

\begin{frame}{实数的 $n$ 进制展开}

\begin{Def}[$n$ 进制展开]
设整数 $n \geqslant 2.$ 对 $x \in [0,1],$ 称
\[
x = \sum_{k=1}^{\infty} \frac{a_k}{n^k}, \quad a_k \in \{0, 1, \ldots, n-1\}
\]
为 $x$ 的 \alert{$n$ 进制展开}.
\end{Def}

\pause

\begin{attnblock}[$n$ 进制展开的非唯一性]
分母为 $n$ 的幂次的有理数有恰好两种展开 (有限展开与无限展开); 其余点展开唯一.
\end{attnblock}

\pause

\begin{eg}
\begin{itemize}
\item $n = 3,$ 三进制:
\[
\dfrac{1}{3} = 0.1_3 = 0.0\overline{2}_3, \quad \dfrac{2}{3} = 0.2_3 = 0.1\overline{2}_3.
\]
\item $n = 10,$ 十进制:
\[
\dfrac{1}{10} = 0.1_{10} = 0.0\overline{9}_{10}, \quad 1 = 1.0_{10} = 0.9\overline{9}_{10}.
\]
\end{itemize}
\end{eg}

\pause
\vspace{0.5em}

一般情况下, 会约定使用有限展开 (或无限展开) 来表示分母为 $n$ 的幂次的有理数, 以保证 $n$ 进制展开的唯一性.
最常用的是约定使用有限展开, 但在构造 Cantor 集时, 约定使用无限展开更为方便.

\end{frame}

%------------------------------------------------

\begin{frame}{实数的 $n$ 进制展开}

\begin{prop}
$[0,1]$ 通过二进制展开 (至多两种) 与 $\{0,1\}^{\mathbb{N}}$ 建立 (几乎处处的) 对等, 故
\[
|[0,1]| = |\{0,1\}^{\mathbb{N}}| = 2^{\aleph_0} = \mathfrak{c}.
\]
\end{prop}

\pause
\vspace{0.5em}

事实上, 记
\begin{align*}
A & = \{ (a_1, a_2, \ldots) : a_k \in \{0,1\} \}, \\
A_0 & = \{ (a_1, a_2, \ldots) \in A : \text{存在 $k$ 使得 $a_k = 0$ 且 $a_j = 1$ 对所有 $j > k$ 都成立} \} \setminus \{(1,1,\ldots)\}.
\end{align*}
那么
\[
A \setminus A_0 \longrightarrow [0, 1]: \quad (a_1, a_2, \ldots) \mapsto \sum_{k=1}^\infty \frac{a_k}{2^k},
\]
即建立了 $A \setminus A_0$ 与 $[0,1]$ 的双射 (二进制的有限表示, 除 $1 = 0.1\overline{1}_2$ 外). 很容易看出, $A_0$ 是可列集, 故 $|A \setminus A_0| = |A| = \mathfrak{c}.$
\qed

\pause
\vspace{0.5em}

\begin{remark}
类似地, $[0,1]$ 可以通过通过一般的 $n$ 进制展开与 $\{0,1,\ldots,n-1\}^{\mathbb{N}}$ 建立 (几乎处处的) 对等.
\end{remark}

\end{frame}

%------------------------------------------------

\section[Cantor 集]{Cantor 三分集}

%------------------------------------------------

\begin{frame}{Cantor 三分集的构造}

从 $[0,1]$ 出发, 逐步移除中间三分之一 (长度为 $1/3^n$) 的开区间:
\begin{align*}
F_1 & = F_{11} \cup F_{12} = \left[0, \tfrac{1}{3}\right] \cup \left[\tfrac{2}{3}, 1\right], &
I_1 & = I_{11} = \left(\tfrac{1}{3}, \tfrac{2}{3}\right), \\[3pt]
F_2 & = \left[0, \tfrac{1}{9}\right] \cup \left[\tfrac{2}{9}, \tfrac{1}{3}\right] \cup \left[\tfrac{2}{3}, \tfrac{7}{9}\right] \cup \left[\tfrac{8}{9}, 1\right], &
I_2 & = \left(\tfrac{1}{9}, \tfrac{2}{9}\right) \cup \left(\tfrac{7}{9}, \tfrac{8}{9}\right), \\[3pt]
F_n & = 2^n \text{ 个闭区间之并,} & I_n &= 2^{n-1} \text{ 个开区间之并,} \\[3pt]
& \vdots & & \vdots
\end{align*}

\pause

令
\[
 P_0 = [0,1] \setminus G_0 = \bigcap_{n=1}^{\infty} F_n \quad \text{(Cantor 三分集)}, \qquad
G_0 = \bigcup_{n=1}^{\infty} I_n \quad \text{(开集)}.
\]

\end{frame}

%------------------------------------------------

\begin{frame}{Cantor 三分集的三进制刻画}

\begin{prop}
\[
P_0 = \left\{ x \in [0,1] : x \text{ 有三进制展开 } x = \sum_{k=1}^{\infty} \frac{a_k}{3^k},\
a_k \in \{0, 2\} \right\}.
\]
\end{prop}

\pause

\begin{remark}
移除 $I_1 = (1/3, 2/3)$ 恰好对应去掉三进制展开首位为 $1$ 的点; 之后每步类似.
端点 $1/3$ 有展开 $0.0\overline{2}_3$ (只含 $0$ 和 $2$), 故 $1/3 \in P_0;$ 类似地 $2/3 \in P_0.$ \alert{注意}, 这里用的是三进制下的无限展开表达.
\end{remark}

\pause

\begin{thm}[$|P_0| = \mathfrak{c}$]
映射 $\varphi: P_0 \to [0,1],$
\[
  \varphi\!\left(\sum_{k=1}^{\infty} \frac{a_k}{3^k}\right) = \sum_{k=1}^{\infty} \frac{a_k/2}{2^k},
  \quad a_k \in \{0, 2\}
\]
是满射 (每个二进制数均可达), 故 $|P_0| \geqslant |[0,1]| = \mathfrak{c},$ 从而 $|P_0| = \mathfrak{c}.$
\end{thm}

\end{frame}

%------------------------------------------------

\begin{frame}{Cantor 三分集的性质}

\begin{prop}[Cantor 三分集 $P_0$ 的基本性质]
\begin{enumerate}
\item \textbf{有界完全集 (闭集)}: $P_0 \subset [0,1];$
$P_0 = \bigcap_{n=1}^\infty F_n$ 为闭集; $P_0$ 无孤立点 ($P_0' = P_0$).
\item \textbf{无处稠密 (无内点)}: $P_0$ 不含任何开区间, $\mathring{P}_0 = \emptyset.$
\item \textbf{测度为零} (预告): 被移除的开集 $G_0$ 的总长度
\[
m G_0 = \sum_{n=1}^{\infty} \frac{2^{n-1}}{3^n} = \frac{1/3}{1 - 2/3} = 1, \quad \text{故 } m P_0 = 0.
\]
(记号 $m$ 表示``长度''概念的推广; 测度的严格定义在第二章给出, 届时将重新推导此式.)
\item \textbf{不可列}: $|P_0| = \mathfrak{c}$ (由满射 $\varphi: P_0 \twoheadrightarrow [0,1]$).
\end{enumerate}
\end{prop}

\pause

\begin{attnblock}
$P_0$ 是测度为零的不可列完全无处稠密集 --- 势与 $\mathbb{R}$ 相同, 但测度为零.
\end{attnblock}

\end{frame}

%------------------------------------------------

\begin{frame}{Cantor 三分集性质的证明}

\startproof[bg=yellow, fg=red]
\textbf{(1) 有界完全集.} $P_0 = \bigcap_{n=1}^\infty F_n$ 是闭集之交, 故为闭集; 显然 $P_0 \subset [0,1]$ 有界.

\pause

\textbf{无孤立点}: 任取 $x \in P_0$, 设其三进制展开 $x = \sum_{k=1}^\infty a_k/3^k$ ($a_k \in \{0, 2\}$). 对每个 $n$, 令
\[
x^{(n)} = \sum_{k \neq n} \frac{a_k}{3^k} + \frac{2 - a_n}{3^n},
\]
即把第 $n$ 位 $a_n$ 换成 $2 - a_n$ ($0 \leftrightarrow 2$). 则 $x^{(n)} \in P_0$, $x^{(n)} \neq x$, 且
\[
|x - x^{(n)}| = \frac{2}{3^n} \to 0 \quad (n \to \infty).
\]
故 $x$ 是 $P_0$ 的聚点. 从而 $P_0' = P_0$, 即 $P_0$ 是完全集.

\pause
\vspace{1em}

\textbf{(2) 无处稠密.} 用反证法: 若 $P^o \neq \emptyset$, 则存在开区间 $(a, b) \subset P_0$. 取 $n$ 充分大使 $3^{-n} < b - a$, 则 $(a, b)$ 必与某个 $F_n$ 的构成区间相交, 但 $F_n$ 的每个构成区间长度为 $3^{-n}$, 且相邻构成区间之间是被移除的开区间 $I_n$, 矛盾. 故 $P^o = \emptyset$.

\end{frame}

%------------------------------------------------

\begin{frame}{Cantor 三分集性质的证明 (续)}

\textbf{(3) 测度为零.} 被移除的开集 $G_0 = \bigcup_{n=1}^\infty I_n$ 的总测度
\[
m G_0 = \sum_{n=1}^\infty \frac{2^{n-1}}{3^n} = \frac{1/3}{1 - 2/3} = 1,
\]
故 $m P_0 = m([0,1]) - m G_0 = 0$.

\pause
\vspace{1em}

\textbf{(4) 不可列.} 映射 $\varphi: P_0 \to [0,1],\; \varphi\!\left(\sum_{k=1}^\infty a_k/3^k\right) = \sum_{k=1}^\infty (a_k/2)/2^k$ 是满射, 故 $|P_0| \geqslant |[0,1]| = \mathfrak{c}$, 而 $P_0 \subset [0,1]$, 因此 $|P_0| = \mathfrak{c}$.
\qed

\end{frame}

%------------------------------------------------

\begin{frame}{再谈不可列性: 区间套法 (另证)}

沿用证明 $[0,1]$ 不可列的思想, 用\alert{闭区间套}可以直接证明 $P_0$ 不可列.

\pause

\startproof[bg=yellow, fg=red](反证)
假设 $P_0$ 可列, 写成点列 $P_0 = \{x_1, x_2, \ldots, x_n, \ldots\}.$ 归纳构造闭区间套: 取 $E_1$ 为 $F_1 = F_{11} \cup F_{12}$ 中不含 $x_1$ 的那个成员; 一般地, 取 $E_n$ 为 $F_{n1} \cup \cdots \cup F_{n2^n}$ 中含于 $E_{n-1}$ 且不含 $x_n$ 的成员 ($x_n$ 至多属于其中之一, 故这样的成员必存在), 则
\[
E_1 \supset E_2 \supset \cdots, \qquad |E_n| = 3^{-n}, \qquad x_n \notin E_n.
\]
\pause
由闭区间套定理, 存在唯一的 $\xi \in \bigcap_{n=1}^\infty E_n.$ 每个 $E_n = F_{n k_n}$ 的两个端点都是被挖去开区间的端点, 因而属于 $P_0,$ 且与 $\xi$ 的距离不超过 $3^{-n} \to 0,$ 故 $\xi$ 是 $P_0$ 的聚点; 又 $P_0$ 是闭集, 故 $\xi \in P_0.$
\pause
于是 $\xi = x_m$ 对某个 $m$ 成立, 但 $\xi \in E_m$ 而 $x_m \notin E_m,$ 矛盾. 故 $P_0$ 不可列.
\qed

\end{frame}

%------------------------------------------------

\graymode
\begin{frame}{补充: 稠密集与疏朗集}

\begin{Def}[稠密集与疏朗集]
设 $E \subset \mathbb{R}.$
\begin{itemize}
\item 若 $\overline{E} = \mathbb{R},$ 称 $E$ 为\alert{稠密集} (dense);
\item 若 $\overline{E}$ 不含任何非空开区间 (即 $\mathring{\overline{E}} = \emptyset$), 称 $E$ 为\alert{疏朗集}, 又称\alert{无处稠密集} (nowhere dense).
\end{itemize}
\end{Def}

\pause

\begin{attnblock}
``$E$ 无内点''与``$E$ 疏朗''不同: $\mathbb{Q}$ 无内点, 但 $\overline{\mathbb{Q}} = \mathbb{R},$ 故 $\mathbb{Q}$ 不是疏朗集.
\end{attnblock}

\pause

\begin{eg}
$P_0$ 是疏朗集: $P_0$ 闭, 且 $\mathscr{C} P_0 = G_0 \cup (-\infty, 0) \cup (1, +\infty)$ 在 $\mathbb{R}$ 中稠密, 故 $\mathring{P}_0 = \emptyset.$
\end{eg}

\end{frame}

\begin{frame}{补充: Baire 纲}

\begin{Def}[第一纲集与第二纲集]
若 $E$ 是\textbf{可数个}无处稠密集之并, 则称 $E$ 为\alert{第一纲集} (first category); 不是第一纲集的集合称为\alert{第二纲集}.
\end{Def}

\pause

\begin{eg}
\begin{itemize}
\item $P_0$ 是无处稠密集, 故是第一纲集.
\item 设 $f(x) = \lim\limits_{k \to \infty} f_k(x)$, 其中 $f_k \in C(\mathbb{R}),$ 则 $f$ 的不连续点集是第一纲集.
\end{itemize}
\end{eg}

\end{frame}
\normalmode

%------------------------------------------------

\section[Cantor 函数]{Cantor 函数与胖 Cantor 集}

%------------------------------------------------

\begin{frame}{Cantor 函数 (恶魔阶梯)}

\begin{Def}[Cantor 函数]
先在 $P_0$ 上定义 $\phi: P_0 \to [0,1]:$ 对 $x = 2\sum_{k=1}^\infty \alpha_k/3^k \in P_0$
($\alpha_k \in \{0,1\}$), 令
\[
  \phi(x) = \sum_{k=1}^{\infty} \frac{\alpha_k}{2^k}.
\]
再扩展为 $[0,1]$ 上的 \alert{Cantor 函数} $\Phi:$
\[
  \Phi: [0,1] \to [0,1], \quad \Phi(x) = \sup\bigl\{\phi(y) : y \in P_0,\ y \leqslant x\bigr\}.
\]
$\Phi$ 在每个构成区间 $(c, d) \subset G_0$ 上取常值 $\phi(c) = \phi(d).$
\end{Def}

\pause

\begin{prop}[$\Phi$ 的性质]
\begin{itemize}
\item $\Phi$ 单调递增, 连续; $\Phi(0) = 0,$ $\Phi(1) = 1.$
\item $\Phi' = 0$ 几乎处处 (在测度为 $1$ 的集 $G_0$ 上导数为 $0$), 但 $\Phi$ 非常数. (\alert{奇异函数})
\item $\Phi$ 连续, 但\alert{不是绝对连续的} (概念见第四章 §6): Newton--Leibniz 公式适用范围的关键反例.
\end{itemize}
\end{prop}

\end{frame}

%------------------------------------------------

\begin{frame}{胖 Cantor 集 (Smith-Volterra-Cantor 集)}

将 Cantor 集构造中每步移除的区间缩短, 可得\alert{正测度}的无处稠密完全集.

\pause

\begin{Def}[胖 Cantor 集]
固定 $0 < \mu < 1.$ 仿照 Cantor 集的构造, 但第 $n$ 步从每个 $2^{n-1}$ 个剩余闭区间各移除中间长度为 $c_n$ 的开区间, 其中
\[
\sum_{n=1}^{\infty} 2^{n-1} c_n = \mu < 1.
\]
所得闭集 $C_\mu = [0,1] \setminus \bigcup_n (\text{移除的开区间})$ 称为\alert{胖 Cantor 集}.
\end{Def}

\pause
\vspace{0.5em}

\begin{eg}
取 $c_n = 2\mu/4^n,$ 则
$\displaystyle\sum_{n=1}^\infty 2^{n-1} \cdot \frac{2\mu}{4^n}
= 2\mu \sum_{n=1}^\infty \frac{1}{2^{n+1}} = 2\mu \cdot \frac{1}{2} = \mu.$
\end{eg}

\end{frame}

%------------------------------------------------

\begin{frame}{胖 Cantor 集 (Smith-Volterra-Cantor 集)}

\begin{prop}
$C_\mu$ 满足: (1) 完全集 (有界、闭、无孤立点); (2) 无处稠密; (3) 测度 $m C_\mu = 1 - \mu > 0.$
\end{prop}

\pause
\vspace{0.5em}

\begin{attnblock}
存在\textbf{正测度的无处稠密完全集} --- 测度与拓扑结构是相互独立的.
\end{attnblock}

\end{frame}

%------------------------------------------------

\section[高维开集]{高维开集与距离函数}

%------------------------------------------------

\begin{frame}{$\mathbb{R}^n$ 中的开集}

$\mathbb{R}^n$ 中点 $x = (x_1, \ldots, x_n)$ 的欧氏范数 $|x| = \bigl(\sum_{i=1}^n x_i^2\bigr)^{1/2}.$

\pause

\begin{Def}[开球与开集]
\begin{itemize}
\item 点 $x$ 的 \alert{$\delta$-邻域 (开球)}: $U(x, \delta) = \{y \in \mathbb{R}^n : |x-y| < \delta\}.$
\item $G \subset \mathbb{R}^n$ 是\alert{开集}, 若每个点都是内点: $\forall ~ x \in G,\ \exists ~ \delta > 0: U(x,\delta) \subset G.$
\end{itemize}
\end{Def}

\pause

\begin{remark}
$\mathbb{R}^n$ 中开集是可列个开球之并 (以有理点为中心、有理数为半径的球构成 $\mathbb{R}^n$ 的可列基). 但一般\textbf{不能写成互不相交开球之并} --- 不同于 $n = 1$ 的情形.
\end{remark}

\pause

\begin{eg}
$\mathbb{R}^2$ 中圆环 $G = \{(x,y) : 1 < x^2+y^2 < 4\}$ 是连通开集, 不能写成互不相交开圆盘之并.
\end{eg}

\end{frame}

%------------------------------------------------

\begin{frame}{矩体与二进矩体}

\begin{Def}[矩体]
在 $\mathbb{R}^n$ 中, 设 $a_i, b_i \in \mathbb{R}$ 且 $a_i < b_i$ ($i = 1, \ldots, n$). 称
\[
\prod_{i=1}^n (a_i, b_i), \qquad \prod_{i=1}^n [a_i, b_i], \qquad \prod_{i=1}^n [a_i, b_i)
\]
分别为\alert{开矩体、闭矩体与半开矩体}.
\end{Def}

\pause

\begin{thm}[构成区间定理的推广]
$\mathbb{R}^n$ 中任一非空开集 $G$ 均可表示为\alert{至多可列个互不相交的半开二进矩体}之并:
\[
G = \bigsqcup_{J \in \mathcal{H}} J, \qquad J = \prod_{i=1}^n \left[ \frac{m_i}{2^k}, \frac{m_i + 1}{2^k} \right), \quad m_i \in \mathbb{Z}.
\]
\end{thm}

\begin{remark}
$G$ 一般不能写成互不相交\textbf{开}矩体之并, 但用半开二进矩体总可以; 这一分解是第二章 $\mathbb{R}^n$ 上测度理论的出发点.
\end{remark}

\end{frame}

%------------------------------------------------

\begin{frame}{矩体与二进矩体 (证明)}

\startproof[bg=yellow, fg=red]
对 $k = 0, 1, 2, \ldots,$ 令 $\Gamma_k$ 为边长 $2^{-k}$、顶点坐标属于 $2^{-k} \mathbb{Z}$ 的半开二进矩体全体. 逐层\alert{剥离}: 记 $\mathcal{H}_k$ 为 $\Gamma_k$ 中含于 $G$、但不含于 $\bigcup_{j<k} \bigcup_{J \in \mathcal{H}_j} J$ 的矩体全体, 并令 $\mathcal{H} = \bigcup_{k \geqslant 0} \mathcal{H}_k.$
\pause
\begin{itemize}
\item \textbf{互不相交}: 同层 $\Gamma_k$ 中两个矩体要么重合要么不交; 不同层之间由剥离保证不交.
\item \textbf{覆盖 $G$}: 任取 $x \in G,$ 由 $G$ 开, 存在 $\delta > 0: U(x, \delta) \subset G.$ 当 $k$ 充分大时, $\Gamma_k$ 中含 $x$ 的矩体直径 $\sqrt{n}/2^k < \delta,$ 故含于 $G,$ 即 $x$ 被某一 $\mathcal{H}_k$ 的成员覆盖.
\item \textbf{至多可列}: 每个非退化矩体必含一个有理点, 不同的矩体含不同的有理点.
\end{itemize}
\qed

\end{frame}

%------------------------------------------------

\begin{frame}{点到集合的距离}

\begin{Def}[点到集合的距离]
设 $\emptyset \neq A \subset \mathbb{R}^n,$ $x \in \mathbb{R}^n.$ 定义\alert{点 $x$ 到集合 $A$ 的距离}:
\[
  d(x, A) = \inf_{y \in A} |x - y|.
\]
\end{Def}

\pause

\begin{prop}[基本性质]
\begin{enumerate}
  \item $d(x, A) = 0 \iff x \in \overline{A}.$
  \item $x \mapsto d(x, A)$ 是 Lipschitz 连续函数 (常数为 $1$):
        \[|d(x, A) - d(z, A)| \leqslant |x - z| \quad (\forall ~ x, z \in \mathbb{R}^n).\]
  \item 若 $A$ 是闭集, 则 $d(x, A) = 0 \iff x \in A.$
\end{enumerate}
\end{prop}

\pause

\startproof[bg=yellow, fg=red](性质 2)
对任意 $y \in A:$ $d(x, A) \leqslant |x-y| \leqslant |x-z| + |z-y|.$ 对 $y$ 取下确界得
$d(x, A) \leqslant |x-z| + d(z, A),$ 即 $d(x,A) - d(z,A) \leqslant |x-z|.$
对称地 $d(z,A) - d(x,A) \leqslant |x-z|.$
\qed

\end{frame}

%------------------------------------------------

\begin{frame}{距离的可达性}

\begin{thm}[点到闭集的距离可达]
设 $\emptyset \neq A \subset \mathbb{R}^n$ 为闭集, 则对任意 $a \in \mathbb{R}^n,$ 存在 $a_0 \in A$ 使得
\[
d(a_0, a) = d(a, A).
\]
\end{thm}

\pause

\startproof[bg=yellow, fg=red](证明)
取充分大的 $\delta > 0$ 使闭球 $\overline{B}(a, \delta) \cap A \neq \emptyset,$ 则 $\overline{B}(a, \delta) \cap A$ 是有界闭集, 且
$d(a, \overline{B}(a, \delta) \cap A) = d(a, A).$
函数 $x \mapsto d(a, x)$ 连续, 在有界闭集 $\overline{B}(a, \delta) \cap A$ 上取到最小值, 即存在 $a_0 \in \overline{B}(a, \delta) \cap A$ 使 $d(a, a_0) = d(a, A).$
\qed

\pause

\begin{cor}[集合间距离的可达性]
设 $A, B \subset \mathbb{R}^n$ 均为非空闭集, 且其中之一有界, 则存在 $a_0 \in A,$ $b_0 \in B$ 使得 $d(a_0, b_0) = d(A, B).$
\end{cor}

\begin{remark}
提示: 不妨设 $A$ 有界, 对连续函数 $x \mapsto d(x, B)$ 在有界闭集 $A$ 上用最值定理, 再对点 $a_0$ 与闭集 $B$ 用上述定理. 详细论证作为\alert{作业习题 23}.
\end{remark}

\end{frame}

%------------------------------------------------

\begin{frame}{集合间的距离}

\begin{Def}[集合间的距离]
设 $A, B \subset \mathbb{R}^n$ 均非空. \alert{集合 $A$ 与 $B$ 的距离}:
\[
  d(A, B) = \inf_{x \in A, ~ y \in B} |x - y| = \inf_{x \in A} d(x, B).
\]
\end{Def}

\pause

\begin{attnblock}[$d(A,B) = 0 \not\Rightarrow A \cap B \neq \emptyset$]
$A = \mathbb{N},$ $B = \bigl\{n + 1/n : n \geqslant 2\bigr\}$ 均是闭集, $A \cap B = \emptyset,$ 但 $d(A, B) = 0.$
\end{attnblock}

\pause

\begin{prop}
若 $A$ 紧 (compact), $B$ 闭, 且 $A \cap B = \emptyset,$ 则 $d(A, B) > 0.$
\end{prop}

\startproof[bg=yellow, fg=red](证明)
函数 $f(x) = d(x, B)$ 在 $\mathbb{R}^n$ 上连续. 由 $A$ 紧, $f$ 在 $A$ 上取到最小值 $m.$ 对每个 $x \in A,$ 由 $x \notin B$ 及 $B$ 闭知 $f(x) = d(x, B) > 0.$ 故 $m > 0,$ 即 $d(A, B) = m > 0.$
\qed

\end{frame}

%------------------------------------------------

\ifIncludeExercise
\begin{frame}{作业}

\begin{block}{教材第一章 §4 习题}
\begin{itemize}
\item \textbf{22.} 设 $F_1, F_2$ 是 $\mathbb{R}^n$ 中的闭集且 $F_1 \cap F_2 = \emptyset,$ 试证存在开集 $G_1, G_2$: $G_1 \cap G_2 = \emptyset,$ $G_i \supset F_i.$ (提示: 考虑函数 $x \mapsto d(x, F_1) - d(x, F_2)$ 的原像.)
\item \textbf{23.} 设 $F_1, F_2$ 为非空闭集, 其中之一有界, 试证存在 $a_1 \in F_1,$ $a_2 \in F_2$ 使 $d(a_1, a_2) = d(F_1, F_2).$ (即上页推论.)
\item \textbf{25.} 试证 $\mathbb{R}^n$ 中每个闭集可表为可列个开集之交, 每个开集可表为可列个闭集之并 (即闭集是 $G_\delta$ 集, 开集是 $F_\sigma$ 集).
\item \textbf{27.} (见教材.)
\end{itemize}
\end{block}

\end{frame}
\fi

%------------------------------------------------

\end{document}
